\documentclass[]{IAC_style_updated}


\begin{document}

\IACpaperyear{2022} % format yyyy
\IACpapernumber{IAC-22,C1,IP,18,x73772} % full paper id
\IAClocation{Paris, France} % used in the header
\IACdate{18-22 September 2022} % used in the header

% Which copyright? If B, put the copyright holder.
\IACcopyrightA{}
%\IACcopyrightB{LateX association}

\title{TudatPy - a versatile open-source astrodynamics software package for space education}

%\IACauthor{Author name}{corresponding affiliation: nr. of corresponding affiliation}{is corresponding author? 0-1}
\IACauthor{Kevin Cowan}{1}{1}
\IACauthor{Filippo Oggioni}{2}{0}
\IACauthor{Geoffrey Garrett}{2}{0}
\IACauthor{Miguel Avillez}{2}{0}
\IACauthor{Joao Encarnacao}{2}{0}
\IACauthor{Marie Fayolle}{2}{0}
\IACauthor{Jonas Hener}{2}{0}
\IACauthor{Marc Naeije}{2}{0}
\IACauthor{Maarten van Nistelrooij}{2}{0}
\IACauthor{Ron Noomen}{2}{0}
\IACauthor{Michael Plumaris}{2}{0}
\IACauthor{Jeremie Gaffarel}{2}{0}
\IACauthor{Sean Cowan}{2}{0}
\IACauthor{Dominic Dirkx}{2}{0}


% Input affiliations here, order is relevant
\IACauthoraffiliation{Space Department, Faculty of Aerospace Engineering, Delft University of Technology, Kluyverweg 1, 2629 HS, Delft, Netherlands \normalfont{E-mail:\ \texttt{k.j.cowan@tudelft.nl}}}
\IACauthoraffiliation{Space Department, Faculty of Aerospace Engineering, Delft University of Technology, Kluyverweg 1, 2629 HS, Delft, Netherlands}

\abstract{
Lorem ipsum dolor sit amet, consectetur adipiscing elit. Integer vestibulum, libero id tempor suscipit, leo mi viverra nulla, non accumsan nulla nunc sed diam. \\
Here's a new paragraph in the abstract.

%
%
%
}


\maketitle
\thispagestyle{fancy} % resets proper header/footer

\fboxrule = 1pt
\fcolorbox{xkcdBlue}{xkcdWhite}{
    \begin{minipage}[]{.9\linewidth}
        Guidance to tudat-team authors: \\
        \textcolor{xkcdRed}{DO NOT EDIT THE MAIN TEX FILE. TO ADD YOUR MATERIAL, PLACE IT IN ONE OF THE RELEVANT TEX FILES, WHICH WILL THEN BE IMPORTED USING AN $\mathtt{\backslash INPUT\lbrace\rbrace}$ STATEMENT.}
    \end{minipage} 
    }

\section{What is TudatPy?}
    Here's an example reference.~\autocite{dirkx2019a} \\
        \lbrack Input required here.\rbrack


    %\lbrack Input required here.\rbrack

\section{Why use TudatPy?}
    \lbrack Target size of text in this section: ½ an A4\rbrack

\lbrack Input required here.\rbrack

Why should someone choose to use TudatPy?
What are the main advantages of using TudatPy (eg. versus other packages → without naming them (?) )?

Advantages / USPs to consider highlighting:
\begin{itemize}
    \item Modularity
    \item Fidelity
    \item Usability, particularly with the Python interface as well as how this allows the use of the vast range of Python tools now available (eg. Numpy, SciPy, Matplotlib, scikit-learn, ...)
    \item Performance
\end{itemize}

\fboxrule = 1pt
\fcolorbox{xkcdBlue}{xkcdWhite}{
    \begin{minipage}[]{.9\linewidth}
        Guidance to tudat-team authors: \\
        To do: 
        \begin{list}{$\square$}
            \item As features are discussed, list published work (or educational settings) in which the various features have been applied.
        \end{list}
    \end{minipage} 
    }
    \lbrack Input required here.\rbrack


\section{Examples of applications}
    %    \lbrack Input required here.\rbrack

The aim for this section, and thus for the case studies (aka examples), is to have all of the cases focused on a single mission/dynamic scenario. The proposed scenario is to focus on a mission from Earth to an asteroid such as Itokawa. NB: A mission to this asteroid is the topic of the example application using PyGMO in the documentation, which can be found here: 
\href{https://docs.tudat.space/en/latest/_src_getting_started/_src_examples/notebooks/pygmo/asteroid_orbit_optimization.html#}{docs.tudat.space}




    \subsection{Analytical state propagation}
        \lbrack Input required here.\rbrack


    Proposal: \\
    Two-part presentation: 
    \begin{enumerate}
        \item Single arc propagation
        \item MGA-DSM
    \end{enumerate}
    \lbrack Input required here.\rbrack

    \subsection{Numerical state propagation}
        \lbrack Input required here.\rbrack



    \subsection{State estimation}
        \lbrack Input required here.\rbrack



    \subsection{Trajectory Optimization}
        \lbrack Input required here.\rbrack


    
%    - Without extensively advertizing for PaGMO \lbrack or/PyGMO\rbrack, mention the integration with PaGMO and the key advantages of its integration, such as its open-source nature, successful track-record, ...

\section{Project setup}
    \lbrack Target size of text in this section: ½ an A4\rbrack

Overview of:

\begin{itemize}
    \item how the Tudat/TudatPy project is set up (eg. tudat vs. tudatpy, conda, github, build+CI pipeline, documentation)
    \item why the project is set up this way (ie. the logic justifying the current setup)
\end{itemize}

\lbrack Input required here.\rbrack


\section{Code setup}
    \lbrack Target size of text in this section: ½ an A4\rbrack

Overview of:

\begin{itemize}
    \item Key points guiding the philosophy of structure the code (in its current guise); eg. why the code was organized into the existing modules;
    \item key points driving the design of the code (in its current guise).
\end{itemize}

\lbrack Input required here.\rbrack


\section{Get started with Tudat}
    A brief indication of where to find it, how to get started ``playing'' with it, how to contribute to development, ...

\fboxrule = 1pt
\fcolorbox{xkcdBlue}{xkcdWhite}{
    \begin{minipage}[]{.9\linewidth}
        Guidance to tudat-team authors: \\
        To do: 
        \begin{list}{$\square$}
            \item Our documentation is rather good, so we do not need to duplicate it extensively here. Instead, the intention is to point to where the reader can get started and briefly highlight how easy it is to do so.
        \end{list}
    \end{minipage} 
    }


\begin{center}
    \begin{minipage}[tc]{30mm}
        This is text in a \LaTeX minipage environment, which could be useful.
    \end{minipage}
\end{center}

\begin{center}
    \begin{figure}[]
    %    \begin{minipage}[c]{.5\linewidth}
                \includegraphics[width=.5\linewidth]{tudelft_logo_black_2022.png}
    %    \end{minipage}
        \caption[]{Caption for a figure (here, an image).}
    \end{figure}
\end{center}


\section*{Acronyms/Abbreviations}
        \lbrack Input required here.\rbrack



\section*{Nomenclature}
    %\lbrack Input required here.\rbrack
\lbrack Use this section if necessary.\rbrack



\section*{Acknowledgements}

\section*{Appendix A (Title)}

\section*{Appendix B (Title)}

%\section*{References}
    \printbibliography
    \label{sec:references}

\end{document}
